\section{Preliminaries}
\label{sec:preliminaries}

This section fixes the probabilistic, type-theoretic, and
finiteness conventions used throughout the chapter, and explains
the most consequential design choices.  We are deliberately
verbose:\ each choice is paired with the alternative we rejected
and a candid statement of what we gained and what we lost.

\subsection{Discrete probability and \TT{PMF}}
\label{sec:prelim:pmf}

Every distribution in the library is a probability mass function
in mathlib's sense:\ a function $\PMF{X}$ assigning to each $x \in
X$ a non-negative extended real number whose total mass equals
one.  The relevant API is monadic — $\PMF{\cdot}$ is the discrete
Giry monad — with \TT{pure}, \TT{bind}, \TT{map}, and a Dirac
embedding $\Dirac_x \in \PMF{X}$.  Expected value of a function
$f : X \to \RR$ under $\mu \in \PMF{X}$ is the (countable) sum
$\sum_{x} \mu(x)\,f(x)$, formalized as
\TT{Math.Probability.expect}.  When $X$ is finite, this is a
finite sum; when $X$ is countably infinite, it is a tsum that
converges by the unit-mass condition together with
non-negativity of $\mu$.

A \emph{kernel} (in the library's sense, \TT{Math.Probability.Kernel})
is a function $\K : A \to \PMF{B}$.  This is not a kernel in the
measure-theoretic sense:\ no $\sigma$-algebra is involved.  It is the
discrete analogue, and is the right object for finite games:\ a
strategy profile maps deterministically to a distribution over
outcomes, and the no-measurability-needed nature of the discrete
case makes it tractable to reason about by equational rewriting.

\paragraph{Why discrete, and what is gained.}  Mechanizing
measure theory in dependent type theory is feasible — mathlib
contains a substantial measure-theoretic library — but the cost is
non-trivial.  A measure on a space $X$ requires a $\sigma$-algebra
on $X$ (a designated collection of subsets closed under
complements and countable unions); functions of interest must be
proved measurable; integration is defined first for
non-negative measurable functions via simple-function
approximation, then extended to integrable real-valued functions;
Fubini–Tonelli, change of variables, and pushforward each require
their own measurability hypotheses.  Even small-looking statements
acquire side conditions — ``measurable function'', ``integrable
function'', ``$\mu$-a.e.\ equal'' — that propagate through every
definition that uses them.  In a finite or countable setting,
these hypotheses are vacuously satisfied (every function from a
finite type to a measurable space is measurable; every function
into $\RR$ from a finite type is integrable), but the
\emph{infrastructure} that produces the side conditions does not
disappear:\ the proof obligations remain to be discharged, and the
notation for distributions is more cumbersome.

In the discrete case the same obligations are vacuous and the
machinery is unnecessary:\ when the underlying type is discrete,
\emph{every} subset is automatically measurable, \emph{every}
function out of it is automatically measurable, and integrals
reduce to ordinary (possibly infinite) sums.  But this only
matters if we adopt the discrete formalism throughout — if the
library is built on the measure-theoretic infrastructure, even
the \emph{statement} of an integral $\int f \, d\mu$ comes with a
well-formedness side condition that $f$ be measurable, and proofs
that should rewrite by ordinary algebra acquire that side
condition at every step.  This is the cost we avoid.

By taking $\PMF{\cdot}$ as the only distribution type, the library
trades expressivity for ergonomics:
\begin{itemize}
\item \emph{Gained.}\ Every distribution-level equation reduces to
a finite-sum identity.  Expected-utility identities follow by
\TT{simp}-style rewriting under the monadic laws.  No
measurability proof obligations arise.  The ``profile $\to$
distribution-over-outcomes'' presentation of strategy spaces is
direct.  Pushforwards are given by \TT{PMF.bind} and
\TT{PMF.map}; products by an explicit \TT{pmfPi}; conditioning by
ordinary support-restricted sums.

\item \emph{Lost.}\ Continuous strategy spaces (a key fragment of
auction theory beyond the canonical models we treat), continuous
type spaces in Bayesian games, and continuous-time stochastic
games are out of scope.  Some classical existence theorems —
notably Glicksberg's existence of a mixed equilibrium for
continuous payoff functions on compact strategy
spaces~\citep{Glicksberg1952}, and existence of equilibria in
Bayesian games with uncountable type spaces~\citep{MilgromWeber1985}
— are not even statable in this idiom.  Information-economic
results that depend on absolute continuity of one
distribution with respect to another are also unavailable.
\end{itemize}

\tradeoff{The library is intentionally a finite-discrete
formalism.  The boundary it places at $\PMF{\cdot}$ is not a
limitation we hope to remove later; it is the design contract.
Downstream users who need continuous payoffs or uncountable type
spaces must take the library as inspiration, not as a base, and
reformulate the relevant definitions over mathlib's measure
theory.}

\subsection{Finiteness}
\label{sec:prelim:finite}

The library uses Lean's \TT{Fintype} typeclass, which carries a
finite enumeration of a type's elements.  ``Finite'' means
\TT{Fintype}, equivalently:\ in bijection with $\{0, \dots, n-1\}$
for some $n$.  The \emph{definitions} of the core objects —
$\PMF{\cdot}$, \TT{KernelGame}, the mixed extension, \TT{pmfPi},
the EU operator, the solution-concept predicates — do not
require \TT{Fintype} on $\Strat_i$ or $\Outcome$:\ strategy and
outcome types may be countably infinite, and the EU operator is
the tsum-based expectation, well-defined on bounded utilities.
Most \emph{existence theorems} carry finite hypotheses where the
proof genuinely needs them — Brouwer's fixed point on a finite
product of simplices, \TT{Finite.exists\_max}, or an enumeration of
pure strategies.  Several of these theorems also have bounded-utility
forms that remove finite outcome carriers while retaining the finite
strategy/tree hypotheses used by the proof.
A parallel family of \TT{\_of\_bounded} lemmas in
\TT{Math.Probability} and \TT{Concepts/MixedExtension} replays
the algebra of expected utility under a $|u| \le C$ hypothesis
without \TT{[Finite \Outcome]}.  Finite hypotheses on the player
type $\Players$ are nearly always present, since the product
$\prod_i \Strat_i$ and the independent product distribution
\TT{pmfPi} are defined over a finite index.  Some declarations
also require \TT{DecidableEq}, e.g.\ everything that uses
\TT{Function.update} for unilateral deviation.

Decidability assumptions enter through \TT{open Classical in},
which turns on the law of excluded middle for the duration of a
single declaration.  This is a benign use of classical reasoning
when it is invoked only to make \TT{Function.update} compute on
abstract index types; on concrete finite types like \TT{Fin n} the
decision is constructive and the classical instance is
definitionally redundant.  We collect these uses in
\S\ref{sec:discussion:noncomputable} and explain why they do
not preclude a future computational refinement.

\subsection{Type-theoretic conventions}
\label{sec:prelim:type-theory}

A reader without a dependent-type-theory background will find
several constructs in the library that look like ordinary set
builders but mean something more precise.  We gloss the ones that
matter:

\begin{itemize}\itemsep2pt
\item \emph{Dependent function types.}  $\prod_i \Strat_i$ in
ordinary mathematics is a product of sets indexed by $\Players$.
In Lean it is the dependent function type \TT{(i : $\iota$) $\to$
Strategy i}, written below as
$\forall i, \Strat_i$.  An element is a function that, for each
$i$, returns a value in the set $\Strat_i$ — exactly a strategy
profile.  The library uses this throughout as the
``strategy-profile'' carrier.

\item \emph{Records / structures.}  A record is a labeled product;
\TT{KernelGame}~$\iota$ bundles four fields into a single value.
Definitional unfolding of structure projections (e.g.\ ``the
\TT{utility} field of $\KG$ is $\util$'') is propositional in Lean
and is freely used in proofs by simp-rewriting.

\item \emph{\TT{Prop} vs.\ \TT{Type}.}  Propositions live in the
sort \TT{Prop} and are proof-irrelevant; data types live in
\TT{Type} (or higher universes) and are proof-relevant.  Solution
concepts are \TT{Prop}-valued; \TT{KernelGame} is \TT{Type}-valued.
The chapter uses ``predicate'' for a \TT{Prop}-valued definition
and ``structure'' for a \TT{Type}-valued one.

\item \emph{Typeclass inference.}  When we write that a definition
``requires \TT{Fintype} on $\Strat_i$'', we mean Lean's typeclass
inference will look up an instance \TT{Fintype $\Strat_i$} from
context.  This is bookkeeping for the reader; the resulting object
is the same whether the instance was synthesized or supplied
explicitly.

\item \emph{Universe levels.}  The library is universe-polymorphic
in $\Players$ and \TT{Outcome} where it has to be, but the chapter
treats them notationally as ordinary sets.  Universe-level
acrobatics arise only at one place — the lifting of solution
concepts across reindexing along an equivalence — and we flag
them when they do.
\end{itemize}

\subsection{The mathlib base}
\label{sec:prelim:mathlib}

The library depends on mathlib~\citep{mathlib2020} for finite-type
infrastructure, the simplex (\TT{stdSimplex} and convex-combination
machinery), product topology, fixed-point theorems (Brouwer in
particular), the discrete Giry monad \TT{PMF}, and a substantial
amount of $\sum$-and-tsum API.  Two non-trivial dependencies
deserve explicit mention:

\begin{itemize}
\item \emph{Brouwer's fixed-point theorem.}  Mixed Nash existence
(\S\ref{sec:theorems:nash-mixed}) is reduced to Brouwer applied to
the product of standard simplices.  Mathlib supplies the
continuous-map and topology infrastructure; the Brouwer theorem
itself is imported from the standalone \TT{fixed-point-theorems-lean4}
package~\citep{harfeFixedPointTheorems}, a Lean~4 mechanization of
Brouwer (and Kakutani) for nonempty compact convex subsets of
finite-dimensional normed spaces.  We use it as a black box.

\item \emph{Schauder's fixed-point theorem.}  The companion file
\TT{Math/SchauderFixedPoint.lean} proves a Schauder fixed-point
principle for nonempty compact convex subsets of real normed
spaces:\ every continuous self-map has a fixed point.  It is
separate from the finite-dimensional product-simplex proof used
for mixed Nash existence.  Its role is analytic infrastructure:
it supplies the fixed-point principle one would need after a
game-theoretic model has supplied the corresponding compactness
and continuity hypotheses.

\item \emph{Discrete Giry monad laws.}  Most equational reasoning
in the library is monadic in \TT{PMF}: bind associativity,
left/right unit, and pushforward laws.  These are mathlib lemmas
we use as black boxes.
\end{itemize}

A small companion file \TT{Math.Probability} adds the few
discrete-probability constructions mathlib does not already
contain (notably \TT{Math.PMFProduct.pmfPi}, the independent
product of an indexed family of $\PMF{\cdot}$'s over a finite
index but otherwise arbitrary per-coordinate types, used in the
mixed extension of a kernel game; and the bounded-utility tsum
identities used to lift expected-utility reasoning past
\TT{[Finite \Outcome]}).

\subsection{What we mean by ``finite game''}
\label{sec:prelim:finite-game}

We will treat a finite game as a kernel game $\KG$ with $\Players$
finite, each $\Strat_i$ finite, and $\Outcome$ finite.  The
language layer adds further structural finiteness — a finite
extensive-form tree, finitely many information sets, a bounded
horizon.  All theorems in the chapter are stated under (some
subset of) these assumptions, made explicit in each statement.
The underlying definitions in the Lean library are stated more
generally — countable $\Strat_i$ and $\Outcome$ are admitted, and
many concept-layer lemmas have bounded-utility variants without a
finite-outcome hypothesis — so a downstream user who needs a
countable-but-not-finite carrier can reuse the definitions and
the algebra; only the major existence theorems are gated by the
finite-game assumptions stated below.
